\documentclass[11pt]{article}
\usepackage{series}
\usepackage{amsmath,amssymb,amsthm}
\usepackage{geometry}
\usepackage{hyperref}
\usepackage{enumitem}

\geometry{margin=1in}


% ---------- metadata ----------
\title{The Modal Triplet Theory Program B1:\\ Gravity as Kinematic Consistency Encoding\\
in the Modal Triplet Theory Program}
\author{Peter Nero}
\date{January 2026}

\begin{document}

\maketitle

\begin{abstract}
We show that gravity arises as a necessary encoding class once coherent kinematics
is defined in the absence of a global reduced description. Building on the
structural core of Modal Triplet Theory and its kinematic formulation, we
demonstrate that stabilization of kinematic persistence across overlapping
descriptions in the presence of circle obstructions requires additional
bookkeeping structure. This structure enforces consistency of admissible
continuation around closed overlap chains and couples all kinematic degrees of
freedom universally.

Gravity is therefore not introduced as a fundamental interaction, dynamical law,
or geometric axiom. It is identified as the unique encoding that resolves
loop-dependent kinematic inconsistency by organizing overlap compatibility.
Geometric notions such as connection, curvature, and metric structure appear as
realizations of this encoding in suitable realization classes, but are not
assumed at the structural level.

This paper establishes gravity as an inevitable response to the coexistence of
local describability, kinematic persistence, and global obstruction, and prepares
the ground for subsequent encoding-class papers on gauge structure and
quantization.
\end{abstract}

\tableofcontents
\newpage

% ============================================================
\section{Introduction and Scope}

The preceding papers in the Modal Triplet Theory Program establish three key
results. First, local reduced descriptions exist only on admissible domains and
cannot be extended globally. Second, kinematics can be defined intrinsically in
terms of admissible continuation across overlapping encodings, without assuming
spacetime or dynamics. Third, failure of global coherence admits an exhaustive
classification into three obstruction types: circle, lens, and nil.

The purpose of the present paper is to analyze the structural consequences of
\emph{circle} obstructions for kinematics. Circle obstructions represent
loop-dependent failure of identity persistence: transporting a coherent
structure along two different closed chains of admissible overlaps yields
inequivalent results. As shown previously, such obstructions are unavoidable in
any system that admits local describability, kinematic persistence, and global
obstruction.

We show that the existence of circle obstructions forces the introduction of an
additional encoding class whose sole role is to enforce consistency of kinematic
continuation across overlapping descriptions. This encoding class is what is
commonly identified as gravity. Crucially, gravity does not appear here as a
force, a field, or a geometric background. It appears as bookkeeping: a structure
that tracks and compensates for loop-dependent inconsistency in kinematic
continuation.

Throughout this paper, we adopt the following guiding principles:
\begin{itemize}
\item Gravity is an \emph{encoding}, not an interaction.
\item Geometry is a \emph{realization} of that encoding, not its definition.
\item Universality of gravitational coupling follows from universality of
kinematic persistence.
\item No dynamical equations are assumed or derived at the structural level.
\end{itemize}

The results of this paper are therefore independent of any specific geometric or
field-theoretic realization. They apply to all realization classes compatible
with the structural core and kinematic framework of Modal Triplet Theory.

\begin{remark}[Position in the series]
This paper depends only on the structural core, coherent kinematics, and the
classification of obstruction types developed earlier in the series. It precedes
papers on gauge structure, quantization, and unified encodings, and should be
read before any realization-specific treatment of gravity or geometry.
\end{remark}

% ============================================================
% Next section will be:
% \section{Circle Obstructions and the Need for Kinematic Bookkeeping}
% ============================================================
\section{Circle Obstructions and the Need for Kinematic Bookkeeping}

In this section we show that coherent kinematics, as defined purely in terms of
admissible continuation across overlapping encodings, is insufficient to ensure
global consistency in the presence of circle obstructions. This insufficiency
forces the introduction of an additional encoding whose sole role is to track
and compensate for loop-dependent kinematic inconsistency.

\subsection{Pure kinematics and path dependence}

Recall that in the kinematic framework, motion is defined as admissible
continuation across chains of overlapping encodings, and worldlines are
equivalence classes of such continuation chains.

In the absence of circle obstructions, admissible continuation is path
independent: transporting a coherent structure along different continuation
chains between the same initial and final encodings yields equivalent results.

However, when a circle obstruction is present, this path independence fails.

\begin{definition}[Kinematic path dependence]
A kinematic system exhibits \emph{path dependence} if there exist admissible
continuation chains
\[
\mathcal E_{\alpha_1} \to \mathcal E_{\alpha_2} \to \cdots \to \mathcal E_{\alpha_n}
\quad \text{and} \quad
\mathcal E_{\alpha_1} \to \mathcal E_{\beta_2} \to \cdots \to \mathcal E_{\alpha_n}
\]
with the same endpoints such that the induced representations of a coherent
structure at $\mathcal E_{\alpha_n}$ are not admissibly equivalent.
\end{definition}

Path dependence is the operational manifestation of a circle obstruction in the
kinematic setting.

\subsection{Breakdown of kinematic identity}

Path dependence has immediate consequences for kinematic identity.

\begin{lemma}
In the presence of a circle obstruction, kinematic identity defined solely by
admissible continuation is ill-defined globally.
\end{lemma}

\begin{proof}
Let two continuation chains with the same endpoints yield inequivalent
representations of a coherent structure. Then there is no encoding-independent
criterion for identifying the two endpoints as representing the same kinematic
state. Therefore identity persistence depends on the path chosen, and global
kinematic identity fails.
\end{proof}

This failure is structural and cannot be repaired by refining the admissible
cover or by choosing different local representatives.

\subsection{Why kinematics alone cannot resolve circle obstructions}

One might attempt to resolve path dependence by modifying the definition of
admissible continuation or by restricting allowed continuation chains. Such
attempts fail.

\begin{lemma}
No restriction of admissible continuation consistent with local describability
can eliminate kinematic path dependence arising from circle obstructions.
\end{lemma}

\begin{proof}
Circle obstructions are invariant under admissible refinement and re-encoding.
Any restriction that removes path dependence on one closed chain will introduce
inconsistency on another, or will eliminate admissible continuation entirely,
collapsing kinematics into nil. Therefore path dependence cannot be resolved
within kinematics alone.
\end{proof}

Thus, pure kinematics is insufficient once circle obstructions exist.

\subsection{Need for additional bookkeeping}

The failure of kinematic identity indicates the need for additional structure.

\begin{remark}
To restore consistency of kinematic persistence, the theory must track how
representations change along different continuation chains and compensate for
path dependence. This requires an additional encoding whose role is not to define
motion, but to enforce consistency of motion.
\end{remark}

This additional encoding must satisfy two requirements:
\begin{enumerate}[label=(\roman*)]
\item it must assign compensating data to continuation chains so that equivalent
endpoints can be identified consistently;
\item it must act universally on all coherent structures, since kinematic
persistence applies universally.
\end{enumerate}

\subsection{Universality of the required encoding}

Because circle obstructions arise from overlap structure rather than from
particular kinds of coherent structure, the required bookkeeping must be
universal.

\begin{lemma}
Any encoding that resolves kinematic path dependence must couple universally to
all coherent structures admitting kinematic persistence.
\end{lemma}

\begin{proof}
Circle obstructions are properties of the encoding atlas and its overlaps, not of
individual coherent structures. Therefore any compensating encoding must apply
uniformly wherever admissible continuation is defined.
\end{proof}

This universality property distinguishes the required encoding from other
encoding classes, such as those addressing lens or nil obstructions.

\subsection{Preview: gravity as bookkeeping}

The observations above motivate the definition of a new encoding class.

\begin{remark}
The encoding required to resolve kinematic path dependence introduced by circle
obstructions is precisely what is identified as gravity. Its function is to
assign compensating data to continuation chains so that kinematic identity is
well-defined despite loop-dependent inconsistency.
\end{remark}

In the next section we define this \emph{kinematic consistency encoding}
abstractly and show that it is uniquely characterized by the requirements
identified here.

% ============================================================
% Next section will be:
% \section{The Kinematic Consistency Encoding}
% ============================================================
\section{The Kinematic Consistency Encoding}

We now introduce the encoding class required to resolve kinematic path dependence
arising from circle obstructions. This encoding is defined abstractly and
structurally, without assuming geometry, fields, or dynamical laws.
\begin{definition}[Path compensation]
Let $\Pi(N)$ denote the path groupoid of the atlas nerve $N$. A \emph{path
compensation} is a functor
\[
C : \Pi(N) \to \mathrm{Aut}(\mathcal E)
\]
assigning to each admissible continuation path an automorphism of encoding data,
such that $C$ is trivial on identity paths and respects path concatenation.
\end{definition}


\subsection{Encoding role and requirements}

From the preceding section, the required encoding must satisfy the following
properties:

\begin{enumerate}[label=(\roman*)]
\item \textbf{Consistency enforcement:} it must compensate for path-dependent
changes in representation induced by admissible continuation around closed
overlap chains;
\item \textbf{Universality:} it must apply uniformly to all coherent structures
admitting kinematic persistence;
\item \textbf{Local definability:} it must be definable on admissible domains and
respect overlap consistency;
\item \textbf{Non-redundancy:} it must not reduce to a lens-type redundancy or to
a restriction of admissible continuation.
\end{enumerate}

Any encoding failing to satisfy these requirements cannot resolve circle
obstructions without collapsing into lens or nil.

\subsection{Definition of the kinematic consistency encoding}

We now define the encoding formally.

\begin{definition}[Kinematic consistency encoding]
A \emph{kinematic consistency encoding} is an admissible encoding class that
assigns to each admissible continuation chain additional compensating data such
that:
\begin{enumerate}[label=(\roman*)]
\item for any two admissible continuation chains with the same endpoints, the
combined effect of continuation and compensation yields admissibly equivalent
representations;
\item compensation is trivial on contractible continuation chains and nontrivial
precisely when a circle obstruction is present;
\item compensation respects overlap consistency and admissible re-encoding.
\end{enumerate}
\end{definition}

The compensating data associated with this encoding do not represent new
coherent structures; they encode consistency conditions on continuation.
\begin{remark}
Lens redundancy acts fiberwise and does not alter transport invariants on closed
paths. By contrast, kinematic consistency compensation acts on continuation paths
and changes loop transport data. The two structures therefore resolve distinct
obstruction types and are not interchangeable.
\end{remark}

\subsection{Resolution of path dependence}

The primary function of the kinematic consistency encoding is to restore
path-independent kinematic identity.

\begin{theorem}[Necessity of kinematic consistency bookkeeping]
\label{thm:necessity-bookkeeping}
Suppose a theory admits local describability, kinematic persistence, and circle
obstructions. If kinematic path dependence is resolved by any admissible
modification of the descriptive framework, then that modification must introduce
a compensating structure assigning data to continuation paths in the atlas nerve
such that:
\begin{enumerate}[label=(\roman*)]
\item compensations compose functorially along path concatenation;
\item compensations are trivial on contractible paths;
\item compensations cancel loop-dependent discrepancies on closed paths.
\end{enumerate}
Any such compensating structure is equivalent, up to admissible re-encoding, to a
kinematic consistency encoding.
\end{theorem}
\begin{corollary}[Resolution of kinematic path dependence]
Kinematic path dependence induced by circle obstructions is resolved if and only
if a kinematic consistency encoding (or an admissibly equivalent encoding) is
introduced.
\end{corollary}


\begin{proof}
(\emph{If}) Given a kinematic consistency encoding, compensating data assigned to
continuation chains ensure that the net effect of transport around any closed
chain is admissibly equivalent to the identity, restoring path-independent
identification of endpoints.

(\emph{Only if}) In the absence of such an encoding, path-dependent inequivalence
of endpoints persists by definition of circle obstruction. No modification of
pure kinematics suffices, as shown previously. Therefore a kinematic consistency
encoding is necessary.
\end{proof}

\subsection{Uniqueness up to admissible equivalence}

We now show that the kinematic consistency encoding is unique in its role.

\begin{theorem}[Uniqueness of the kinematic consistency encoding]
Any admissible encoding that resolves kinematic path dependence induced by circle
obstructions is equivalent, up to admissible re-encoding, to a kinematic
consistency encoding.
\end{theorem}

\begin{proof}
Let $\mathcal G$ be an admissible encoding resolving path dependence. By
universality, $\mathcal G$ assigns compensating data to all continuation chains.
By local definability and overlap consistency, this data must compose consistently
on overlaps. By non-redundancy, it cannot be absorbed into local representative
choice. Therefore $\mathcal G$ satisfies the defining properties of a kinematic
consistency encoding and differs from any other such encoding only by admissible
re-encoding.
\end{proof}

\subsection{Gravity as encoding}

We are now in a position to identify gravity within the MTT framework.

\begin{remark}[Gravity as encoding]
The kinematic consistency encoding defined above is what is conventionally
identified as gravity. In Modal Triplet Theory, gravity is not a force or a
dynamical field, but the unique encoding that enforces consistency of kinematic
persistence in the presence of circle obstructions.
\end{remark}

This identification is structural and does not depend on any particular
realization.

\subsection{No dynamical assumptions}

We emphasize that no equations of motion or dynamical laws have been introduced.

\begin{remark}
The kinematic consistency encoding specifies \emph{what} bookkeeping is required
to resolve path dependence, not \emph{how} that bookkeeping evolves. Dynamical
equations, when present, arise only in specific realization classes and are not
part of the encoding definition.
\end{remark}

\subsection{Preview: geometric realizations}

In the next section we show that, in realization classes admitting smooth
structure, the kinematic consistency encoding naturally takes the form of
geometric data such as connections and curvature. These structures will be shown
to be realizations of the encoding defined here, not its definition.

% ============================================================
% Next section will be:
% \section{Geometric Realizations of the Kinematic Consistency Encoding}
% ============================================================
\section{Geometric Realizations of the Kinematic Consistency Encoding}

In this section we show how the kinematic consistency encoding introduced above
is realized in classes of models admitting smooth structure. Geometry appears
here not as a foundational assumption, but as the minimal mathematical framework
capable of implementing the compensating data required to resolve kinematic path
dependence.

\subsection{Realization assumptions}

We explicitly restrict attention to realization classes in which smooth
structures are available.

\begin{assumption}[Smooth realization class]
\label{ass:smooth-realization}
Assume that admissible encodings admit realizations on smooth manifolds such
that:
\begin{enumerate}[label=(\roman*)]
\item admissible domains are represented by smooth coordinate charts;
\item admissible re-encodings on overlaps are smooth transition maps;
\item admissible continuation chains admit smooth refinement.
\end{enumerate}
\end{assumption}

This assumption does not constrain the structural results of the previous
sections; it specifies only the realization class in which geometric objects can
be meaningfully defined.

\subsection{Transition data and bundle structure}

We now show that overlap consistency and compensating data naturally organize
into bundle structure.

\begin{theorem}[Bundle realization of consistency encoding]
\label{thm:bundle-realization}
Under Assumption~\ref{ass:smooth-realization}, the transition data associated
with admissible re-encodings and kinematic consistency compensation define a
principal bundle over each admissible domain, unique up to admissible
equivalence.
\end{theorem}

\begin{remark}[Cocycle versus curvature]
Overlap consistency ensures that transition maps satisfy the Čech cocycle
condition and therefore define a principal bundle. Circle obstructions do not
correspond to failure of the cocycle condition, but to failure of flatness of the
associated transport structure. The role of the kinematic consistency encoding
is therefore not to repair the bundle cocycle, but to introduce a connection
whose curvature encodes the circle obstruction as nontrivial holonomy.
\end{remark}


\subsection{Connections as realization of compensating data}

We now identify the geometric object that realizes the compensating data of the
kinematic consistency encoding.

\begin{theorem}[Connection realization]
\label{thm:connection-realization}
In a smooth realization class, the compensating data of the kinematic consistency
encoding is realized by a connection on the principal bundle defined in
Theorem~\ref{thm:bundle-realization}.
\end{theorem}

\begin{proof}
Kinematic compensation assigns data to infinitesimal continuation segments such
that path-dependent discrepancies cancel on closed loops. In a smooth setting,
this is equivalent to specifying parallel transport rules along infinitesimal
paths.

The unique structure implementing infinitesimal parallel transport compatible
with bundle transition data is a connection. Therefore the compensating data is
realized by a connection on the principal bundle.
\end{proof}

\subsection{Curvature as circle obstruction}

We now show that curvature is the geometric manifestation of circle obstruction.

\begin{theorem}[Curvature realization of circle]
\label{thm:curvature-circle}
In a smooth realization class, the curvature of the connection associated with
the kinematic consistency encoding vanishes if and only if no circle obstruction
is present.
\end{theorem}

\begin{proof}
If the curvature vanishes, parallel transport is path independent on contractible
loops, and no circle obstruction exists. Conversely, if a circle obstruction is
present, path-dependent transport persists, which is detected infinitesimally by
nonvanishing curvature. This equivalence is standard in the theory of connections
and holonomy.
\end{proof}

Thus, curvature is not an additional postulate; it is the realization of the
structural circle invariant.

\subsection{Metric structure as secondary}

We emphasize that metric structure is not required at the encoding level.

\begin{remark}
The kinematic consistency encoding requires only parallel transport and holonomy.
Metric structure, when present, serves to compare magnitudes or define action
functionals, but is not required to define or resolve circle obstructions.
\end{remark}

Metrics therefore appear only in further specialization of realization classes,
not as part of the definition of gravity.

\subsection{Universality of geometric coupling}

We now recover the universal coupling property of gravity.

\begin{theorem}[Universality of coupling]
\label{thm:universal-coupling}
In any smooth realization of the kinematic consistency encoding, the associated
geometric structures couple universally to all coherent structures admitting
kinematic persistence.
\end{theorem}

\begin{proof}
By definition, the kinematic consistency encoding applies to all admissible
continuation chains. Therefore the associated connection and curvature affect
parallel transport of all coherent structures equally. No encoding-relative
charge or selection can screen this coupling.
\end{proof}

This universality is structural rather than dynamical.

\subsection{Summary}

In smooth realization classes, the kinematic consistency encoding is realized as
bundle geometry equipped with a connection whose curvature encodes circle
obstructions. Geometry therefore appears as bookkeeping required to stabilize
kinematic identity under continuation, not as a foundational input.

% ============================================================
% Next section will be:
% \section{Relation to Causal Structure and Horizons}
% ============================================================
\section{Relation to Causal Structure and Horizons}

In this section we connect the kinematic consistency encoding to the causal
structure derived in coherent kinematics. We show that gravity, understood as
bookkeeping of kinematic consistency, is required to stabilize causal cones,
null boundaries, and horizons under refinement.

\subsection{Causal cones under admissible continuation}

Recall from the kinematic framework that admissible continuation induces a causal
partial order on encodings and worldlines, with associated causal cones.

In the absence of circle obstructions, causal cones are path independent and
stable under refinement. However, when circle obstructions are present, causal
relations themselves become path dependent.

\begin{lemma}
In the presence of a circle obstruction, causal precedence defined purely by
admissible continuation is path dependent.
\end{lemma}

\begin{proof}
A circle obstruction implies the existence of two admissible continuation chains
with the same endpoints but inequivalent induced representations. Since causal
precedence is defined in terms of admissible continuation, its evaluation depends
on the continuation path chosen. Therefore causal relations are not invariant
under refinement. By the definition of worldline equivalence in the coherent kinematics framework,
causal precedence is defined in terms of admissible continuation yielding
equivalent endpoint representations, not merely the existence of any path.

\end{proof}

This instability is incompatible with a consistent notion of causal structure.

\subsection{Stabilization of causal structure}

We now show that the kinematic consistency encoding stabilizes causal relations.

\begin{theorem}[Causal stabilization]
\label{thm:causal-stabilization}
In the presence of circle obstructions, causal cones and null boundaries are
well-defined and path independent if and only if the kinematic consistency
encoding is introduced.
\end{theorem}

\begin{proof}
(\emph{If}) The compensating data provided by the kinematic consistency encoding
ensures that continuation along different paths yields admissibly equivalent
results. Therefore causal precedence and null/timelike classification are
independent of the chosen continuation chain.

(\emph{Only if}) Without compensating data, path-dependent continuation persists
by definition of circle obstruction, and causal relations depend on the path
chosen. Thus causal structure cannot be stabilized.
\end{proof}

\subsection{Null boundaries and geometric realization}

Null boundaries identified in coherent kinematics correspond to minimal
admissible continuation.

\begin{remark}
In smooth realizations, stabilization of null boundaries by the kinematic
consistency encoding corresponds to the emergence of null hypersurfaces with
well-defined geometric properties. This correspondence is realizational rather
than definitional.
\end{remark}

Thus, structures commonly identified as light cones arise as stabilized
kinematic boundaries.

\subsection{Horizons and termination}

We now examine the relation between gravity and kinematic horizons.

\begin{theorem}[Horizon stabilization]
\label{thm:horizon-stabilization}
In the presence of circle obstructions, the location and persistence of
kinematic horizons are invariant under refinement if and only if the kinematic
consistency encoding is present.
\end{theorem}

\begin{proof}
Horizons are boundaries of admissible continuation. Without compensating data,
the position of such boundaries depends on the continuation path and refinement
choice. With the kinematic consistency encoding, admissible continuation is
path independent up to equivalence, and horizon boundaries are therefore stable
features of the kinematic structure.
\end{proof}

This shows that gravity stabilizes not only motion but also the termination of
motion.

\subsection{Irreversibility and gravitational bookkeeping}

We revisit irreversibility in the presence of gravity-as-encoding.

\begin{remark}
Irreversibility derived kinematically remains irreversible under the introduction
of the kinematic consistency encoding. Gravity stabilizes irreversible structure
but does not eliminate it. This explains why gravitational dynamics respects but
does not reverse the arrow of kinematic irreversibility.
\end{remark}

\subsection{Preview: dynamical realizations}

The present paper has treated gravity purely as an encoding enforcing kinematic
consistency. In subsequent realization papers, one may introduce action
principles or dynamical equations governing the evolution of the geometric data.
Such dynamics are not required to define gravity as an encoding and are therefore
logically downstream of the present analysis.

% ============================================================
% Next section will be:
% \section{Summary and Outlook}
% ============================================================
\section{Summary and Outlook}

In this paper we have shown that gravity arises as a necessary encoding class
once coherent kinematics is defined in the absence of global reduced description.
Starting from the structural core of Modal Triplet Theory and its kinematic
formulation, we demonstrated that circle obstructions introduce unavoidable
path dependence in admissible continuation. Pure kinematics is insufficient to
resolve this inconsistency.

We proved that resolution of kinematic path dependence requires the introduction
of an additional encoding whose sole role is to enforce consistency of identity
persistence across overlapping descriptions. This encoding is uniquely
characterized by universality, non-redundancy, and local definability. Within
the MTT framework, this encoding is identified as gravity.

Crucially, gravity is not introduced here as a force, a field, or a geometric
axiom. It is defined as bookkeeping: an encoding that compensates for
loop-dependent kinematic inconsistency. In realization classes admitting smooth
structure, this bookkeeping is realized as bundle geometry with connection and
curvature, but such geometric structures are consequences of the encoding rather
than its definition.

We further showed that gravity stabilizes causal structure derived from coherent
kinematics. Causal cones, null boundaries, and horizons become invariant under
refinement only when the kinematic consistency encoding is present. Gravity
therefore couples universally to all coherent structures admitting kinematic
persistence, not because of an interaction postulate, but because kinematics
itself is universal.

The results of this paper complete the identification of gravity within the
Modal Triplet Theory Program. Together with the prior classification of circle,
lens, and nil obstructions, they establish gravity as the unique response to
loop-dependent failure of global coherence.

Subsequent papers in the series develop the remaining encoding classes forced by
structural obstructions. Gauge structure is treated as a redundancy encoding
addressing lens obstructions, while quantization is treated as a discrete
constraint encoding associated with nil and topological termination. Unified
encoding frameworks arise when multiple obstruction types are resolved
simultaneously. Separate realization papers introduce dynamics and explicit
geometric models without altering the structural conclusions established here.

In this way, gravity appears not as a fundamental ingredient of the theory, but
as an inevitable encoding response to the coexistence of local describability,
kinematic persistence, and global obstruction.

% ============================================================
% End of Paper B3

\end{document}
